% 02_problem_characterization.tex — Problem Characterization
\section{Problem Characterization}
\label{sec:problem}

\subsection{Structural Invariants}

We define the class of \emph{coupled decision problems with heterogeneous knowledge} through three invariants that hold across domains (Table~\ref{tab:domains}).

\paragraph{Invariant 1: Lever Coupling.}
Let $\mathcal{L} = \{l_1, \ldots, l_n\}$ denote the set of decision levers.
The optimization landscape exhibits coupling when $\frac{\partial^2 f}{\partial l_i \partial l_j} \neq 0$ for at least one pair $(l_i, l_j)$.
In the Revenue Growth Management domain, for instance, pricing and promotion decisions are coupled because promotional discounts interact with base-price elasticity---a 30\% promotion on a high-base-price SKU produces different revenue effects than the same discount on a low-base-price SKU.

\paragraph{Invariant 2: Knowledge Heterogeneity.}
The information required for decision-making is distributed across $K$ knowledge sources $\{S_1, \ldots, S_K\}$ with incompatible representational schemas.
Quantitative data provides statistical signals, policy documents encode hard and soft constraints, and expert judgment contributes qualitative directional beliefs.
Naive concatenation destroys the epistemic metadata that distinguishes a high-confidence statistical estimate from an intuition-based expert belief.

\paragraph{Invariant 3: Cognitive Assembly Bottleneck.}
Even when all knowledge sources are available, the cognitive load of jointly reasoning across multiple coupled levers and heterogeneous sources exceeds the capacity of unaided human or single-pass algorithmic analysis.
The combinatorial explosion of lever interactions ($\mathcal{O}(2^{|\mathcal{L}|})$ pairwise) combined with the need to reconcile conflicting signals across knowledge types creates a synthesis challenge that requires structured decomposition.

\subsection{Insufficiency of Isolation}

\begin{proposition}
Addressing any single invariant in isolation produces solutions that are necessarily incomplete.
\end{proposition}

\textit{Argument.}
Without coupling awareness, the optimizer treats each lever independently, missing interaction effects such as the price-promotion trap (GT-1) and the assortment-distribution synergy (GT-2).
Without heterogeneous encoding, cross-source signals---where the correct action requires synthesizing quantitative data, policy constraints, and expert judgment (GT-4)---cannot be detected.
Without progressive assembly, the cognitive bottleneck prevents systematic exploration of the ${\sim}10^{18}$ combinatorial space, resulting in ad-hoc solutions with constraint violations.
Our ablation study (\S\ref{sec:results}) provides empirical evidence for this proposition.

\begin{table}[ht]
\centering
\caption{Decision domains sharing the three structural invariants.}
\label{tab:domains}
\small
\begin{tabular}{lcccp{4cm}}
\toprule
\textbf{Domain} & \textbf{Lever Coupling} & \textbf{Knowledge Het.} & \textbf{Cognitive Asm.} & \textbf{Example Levers} \\
\midrule
Revenue Growth Mgmt & \checkmark & \checkmark & \checkmark & Pricing, promotion, assortment, distribution, pack-price \\
Precision Medicine & \checkmark & \checkmark & \checkmark & Drug, dose, timing \\
Materials Discovery & \checkmark & \checkmark & \checkmark & Composition, processing, structure \\
Supply Chain Design & \checkmark & \checkmark & \checkmark & Sourcing, inventory, routing \\
Financial Risk Mgmt & \checkmark & \checkmark & \checkmark & Hedging, allocation, timing \\
\bottomrule
\end{tabular}
\end{table}
